% !PRINTED none
% [part-title page: no running-head folio visible on this scan; opens Dritter Abschnitt and § 27]
\begin{center}
\gesperrt{Dritter Abschnitt.}

\textbf{\gesperrt{Transformation der Theta-Funktionen.}}
\end{center}

\noindent\hrulefill

\begin{center}
\textbf{§ 27. \gesperrt{Das Transformationsprinzip.}}
\end{center}

Wir kehren zurück zu den in § 17 gegebenen Definitionsgleichungen
der $T$-Funktionen $m$ter Ordnung und versehen darin
aus einem gleich ersichtlichen Grunde die Buchstaben $\omega$, $a$, $b$, $m$
mit Akzenten, so daß diese Gleichungen lauten:
\weq{1}{\begin{gathered}
T(u+\omega'_1) = e^{-\pi i [a'_1(2u+\omega'_1)+b'_1]} T(u), \\
T(u+\omega'_2) = e^{-\pi i a'_2 (2u+\omega'_2)+b'_2} T(u). % [sic] Printed without brackets around the full exponent.
\end{gathered}}
\weq{2}{a'_2\omega'_1 - a'_1\omega'_2 = m'.}
Sind $a$, $c$ irgend welche ganze (positive oder negative) Zahlen,
so ergibt sich, wenn man in § 17, (16) $n_1$, $n_2$ durch $-c$, $a$ ersetzt:
\weq{3}{\begin{gathered}
T(u - c\omega'_1 + a\omega'_2) \\
= e^{-\pi i(-ca'_1+aa'_2)(2u-c\omega'_1+a\omega'_2) - \pi i(-cb'_1+ab'_2-m'ac)} T(u),
\end{gathered}}
eine Gleichung, die auch aus einer der Gleichungen (1) hervorgeht,
wenn man darin $a'$, $b'$, $\omega'$ durch
\[
-ca'_1+aa'_2,\quad -cb'_1+ab'_2-m'ac,\quad -c\omega'_1+a\omega'_2
\]
ersetzt.

Hierin ist das Prinzip der Transformation der $T$-Funktionen
enthalten.

Es seien $b$, $\partial$ zwei andere ganze Zahlen, für welche die
Determinante
\weq{4}{n = a\partial - bc}
einen \gesperrt{positiven} Wert hat. Wir setzen
\weq{5}{\begin{gathered}
\omega_1 = +\partial\omega'_1 - b\omega'_2, \\
\omega_2 = -c\omega'_1 + a\omega'_2,
\end{gathered}}
und folglich
\weq{6}{\begin{gathered}
n\omega'_1 = a\omega_1 + b\omega_2 \\
n\omega'_2 = c\omega_1 + \partial\omega_2.
\end{gathered}}
